\documentclass[11pt,a4paper]{article}
\usepackage[a4paper,left=28mm,right=38mm,top=25mm,bottom=25mm,marginparwidth=30mm]{geometry}
\usepackage{amsmath,amssymb,amsthm,mathtools}
\usepackage{booktabs,tabularx,array,enumitem}
\usepackage[T1]{fontenc}
\usepackage{lmodern}
\usepackage{microtype}
\usepackage{xcolor}
\usepackage{setspace}
\usepackage{fancyhdr}
\usepackage[hidelinks]{hyperref}
\usepackage{xurl}
\setstretch{1.18}
\setlength{\parskip}{4pt}
\setlength{\parindent}{0pt}

\definecolor{boxgray}{gray}{0.93}
\newcommand{\statusbox}[1]{\par\smallskip\noindent\colorbox{boxgray}{\parbox{\dimexpr\linewidth-2\fboxsep}{\small #1}}\par\smallskip}

\theoremstyle{plain}
\newtheorem{lemma}{Lemma}
\newtheorem*{lemmaSone}{Lemma S1$'$}
\newtheorem*{lemmaB}{Lemma B}
\newtheorem*{corollary}{Corollary (Gate 1A FAIL)}
\theoremstyle{definition}
\newtheorem{definition}{Definition}
\theoremstyle{remark}
\newtheorem*{remark}{Remark}

\newcommand{\eps}{\epsilon}
\newcommand{\R}{\mathbb{R}}
\newcommand{\Pc}{P_{c}}
\newcommand{\norm}[1]{\lVert #1\rVert}

\pagestyle{fancy}
\fancyhf{}
\lhead{\small sense-correct-abstain}
\rhead{\small Lemma S1$'$ and Lemma B --- draft for review}
\cfoot{\small \thepage}

\begin{document}

\begin{center}
{\LARGE\bfseries Lemma S1$'$ and Lemma B}\\[4pt]
{\large Draft proofs for owner review}\\[8pt]
{\small Source: \texttt{docs/errata/E006\_audit\_of\_9201ba9.md}, \S2--\S3, commit \texttt{4867fe3} (2026-11-21)}
\end{center}

\statusbox{\textbf{Status.} Both lemmas are \emph{proved (draft, owner review pending)}. Numerical checks are deterministic
evaluations of frozen circuits and are \emph{not} proofs. Scope: the synthetic compiled-angle E2 model
(no physical carrier established), the frozen Gate 1A designs (Lemma S1$'$), and the frozen
$p=1$ QAOA MaxCut circuits (Lemma B). Nothing here is a hardware claim. This document only
typesets E006; if the two disagree, the repository file governs.}

\tableofcontents
\vspace{6pt}

%=====================================================================
\section{Three different notions}

Let $\{P_{\eps,\eta}\}$ be a statistical model with parameter of interest $\eps$ and nuisance $\eta$,
and let $(\eps_0,\eta_0)$ be the true point.

\begin{definition}[Zero local efficient information]
The score for $\eps$ at $(\eps_0,\eta_0)$ lies in the closed span of the nuisance tangent set,
so its projection onto the orthocomplement vanishes. This is a \emph{first-order} statement.
It does \emph{not} by itself imply non-identification: a parameter can be identified at higher
order while having zero Fisher information.
\end{definition}

\begin{definition}[Local non-identification]
Every neighbourhood of $\eps_0$ contains some $\eps'\neq\eps_0$ for which there is an $\eta'$ with
$P_{\eps',\eta'}=P_{\eps_0,\eta_0}$ \emph{exactly}.
\end{definition}

\begin{definition}[Global non-identification]
For every $\eps'$ in the parameter set there is an $\eta'$ with $P_{\eps',\eta'}=P_{\eps_0,\eta_0}$.
\end{definition}

\textbf{Rule used throughout.} Every non-identification claim below is proved by an
\emph{explicit construction of observationally equivalent parameter points}, never from
Fisher information alone.

%=====================================================================
\section{\texorpdfstring{Lemma S1$'$ (RC labels at $\phi=0$)}{Lemma S1' (RC labels at phi=0)}}

\subsection{Setting}

\begin{itemize}[leftmargin=*,itemsep=2pt]
\item \textbf{Data.} For each candidate circuit $c$ (tagged logical angles $\theta_{c,\ell}$, $\ell=1,\dots,L_c$)
and each Pauli label pattern $t$ with sign vector $s(t)\in\{\pm1\}^{L_c}$, i.i.d.\ draws from an outcome
law $P_{c,t}$ on the $2^n$ bitstrings. Labels are drawn uniformly and recorded.
\item \textbf{Primary synthetic model.} Each physical tagged rotation with commanded angle
$a=s_\ell(t)\theta_{c,\ell}$ is implemented as $R_Z\big((1+\eps)a+\phi\big)$, and Z labels are
error-free virtual frames.
\item \textbf{Frame covariance (proved at M2).} Undoing the label gives logical angles
\[
u_{c,\ell}(t) = (1+\eps)\,\theta_{c,\ell} + s_\ell(t)\,\phi .
\]
\item \textbf{Key fact.} At $\phi=0$, $u_c(t)=(1+\eps)\theta_c$ for every $t$, so
\[
P_{c,t}=P_c\quad\text{for all label patterns } t.
\]
That is, the label-resolved laws are label invariant (verified to $0$ in the frozen run).
\end{itemize}

\subsection{Three nuisance parameterizations}

The frozen specification (\S5.3) calls its nuisance class ``saturated'' but describes it through
``local response directions'', so it admits more than one reading. The lemma is therefore stated
for three explicit parameterizations.

\medskip
\noindent\begin{tabularx}{\linewidth}{@{}l>{\raggedright\arraybackslash}X>{\raggedright\arraybackslash}X@{}}
\toprule
Name & Law of $(c,t)$ & Nuisance \\
\midrule
\textbf{N-free} & $P_{c,t}=R_c\big(u_c(t)\big)$ & for each $c$, an arbitrary map $R_c$ from angle vectors to strictly positive distributions (the exact quantum kernel is one member) \\[3pt]
\textbf{N-loc} & $P_{c,t}=q_c+\sum_\ell\big(\eps\theta_{c,\ell}+s_\ell(t)\phi\big)D_{c,\ell}$ & $q_c$ in the open simplex; $D_{c,\ell}$ sum-zero vectors, free; positivity of $P_{c,t}$ required \\[3pt]
\textbf{N-q} & as N-loc & only $q_c$ unknown; $D_{c,\ell}$ fixed or known \\
\bottomrule
\end{tabularx}

\medskip
N-q is \emph{not} a reading of the frozen class. It is the weakest assumption under which the
Gate 1A failure still follows, which is why it is included.

Throughout, write $v_c := \sum_\ell \theta_{c,\ell} D_{c,\ell}$ (a sum-zero vector). At $\phi=0$,
N-loc and N-q give $P_c = q_c + \eps\, v_c$.

\subsection{Statement}

\begin{lemmaSone}
At every parameter point with $\phi=0$:
\begin{enumerate}[label=\textup{(\arabic*)},itemsep=3pt]
\item \textbf{Local non-identification} under N-q, and hence under N-loc.
\item \textbf{Global non-identification} under N-loc.
\item \textbf{Global non-identification} under N-free (over all $\eps'$ with $1+\eps'\neq0$).
\item \textbf{Zero local efficient information} for $\eps$ under N-q, N-loc and N-free, assuming the law is
differentiable in $\eps$ and strictly positive (true for the frozen kernel).
\end{enumerate}
\end{lemmaSone}

\subsection{Proofs}

\paragraph{Proof of (1): local construction under N-q.}
\begin{enumerate}[label=\textbf{Step \arabic*.},leftmargin=*,itemsep=2pt]
\item Let $\delta := \min_c \min_b q_c(b) \,/\, \max_c \norm{v_c}_\infty$. If every $v_c=0$, take $\delta=\infty$.
Then $\delta>0$, because each $q_c$ lies in the \emph{open} simplex.
\item Fix any $\eps'$ with $|\eps'-\eps_0|<\delta$. Define $q'_c := q_c + (\eps_0-\eps')\,v_c$ and keep $D$ unchanged.
\item $q'_c$ sums to one, because $v_c$ sums to zero. It is strictly positive, because
$|(\eps_0-\eps')v_c(b)| < \delta\norm{v_c}_\infty \le q_c(b)$.
\item For every $t$: $\;q'_c+\eps' v_c = q_c + \eps_0 v_c$. Hence $P'_{c,t}=P_{c,t}$ for all $c,t$.
\end{enumerate}
Every $\eps'$ in the open interval $(\eps_0-\delta,\eps_0+\delta)$ is therefore observationally equivalent to $\eps_0$. \qed

\paragraph{Proof of (2): global construction under N-loc.}
\begin{enumerate}[label=\textbf{Step \arabic*.},leftmargin=*,itemsep=2pt]
\item If $\eps'\neq0$: set $D'_{c,\ell}:=(\eps_0/\eps')D_{c,\ell}$ and $q'_c:=q_c$. Then $\eps' v'_c=\eps_0 v_c$,
so $q'_c+\eps' v'_c = P_c$ for every $t$.
\item If $\eps'=0$: set $q'_c := P_c$ (strictly positive) and $D':=0$. Again $P'_{c,t}=P_c$.
\item Positivity is automatic in both cases, because the resulting laws equal the true, strictly positive laws.
\end{enumerate}
So every real $\eps'$ is observationally equivalent to $\eps_0$. \qed

\paragraph{Proof of (3): global construction under N-free.}
\begin{enumerate}[label=\textbf{Step \arabic*.},leftmargin=*,itemsep=2pt]
\item At $\phi=0$ the data reveal $R_c$ only at the single point $(1+\eps_0)\theta_c$.
\item For any $\eps'$ with $1+\eps'\neq0$, define $R'_c$ to equal $R_c\big((1+\eps_0)\theta_c\big)$ at the
point $(1+\eps')\theta_c$, and to be any strictly positive distribution elsewhere.
\item Then $P'_{c,t}=R'_c\big((1+\eps')\theta_c\big)=R_c\big((1+\eps_0)\theta_c\big)=P_{c,t}$ for every $t$. \qed
\end{enumerate}

\begin{remark}
No global rescaling-closure property of the response class is used. That premise appeared in an
earlier draft and was rejected at M2. Only one point per candidate is ever observed at $\phi=0$.
\end{remark}

\paragraph{Proof of (4): zero efficient information.}
\begin{enumerate}[label=\textbf{Step \arabic*.},leftmargin=*,itemsep=2pt]
\item At $\phi=0$, the derivative of $P_{c,t}$ with respect to $\eps$ is $v_c$ (N-loc, N-q) or
$\theta_c\cdot\nabla R_c$ (N-free). In both cases it is \emph{the same for every label} $t$ and sums
to zero over outcomes.
\item A label-independent, sum-zero perturbation of the per-candidate distribution is a nuisance
tangent direction in every one of the three classes (free $q_c$, or free $R_c$).
\item The $\eps$ score therefore lies in the nuisance tangent span, and its efficient projection is zero. \qed
\end{enumerate}

\subsection{\texorpdfstring{From Lemma S1$'$ to the Gate 1A verdict}{From Lemma S1' to the Gate 1A verdict}}

\begin{corollary}
Under every reading of the frozen nuisance class, Gate 1A fails its frozen rule.
\end{corollary}
\begin{proof}
\begin{enumerate}[label=\textbf{Step \arabic*.},leftmargin=*,itemsep=2pt]
\item The frozen Gate 1A estimator is the profile maximum-likelihood estimator under the I3 class
(specification \S7). It is applied to \emph{exact probabilities} at every grid point and must reach
``the same identified solution'' from all starts, within $10^{-8}$.
\item The frozen $9\times9$ grid contains nine points with $\phi=0$. At each of them, item (1)
gives an open interval of $\eps'$ whose laws equal the truth exactly. The profile likelihood is
therefore maximal on that whole interval, and the maximiser is not unique.
\item Specification \S7 states that a flat profile or nonunique solution ``is a failure'';
specification \S9 requires recovery ``at every required grid point''.
\item Hence Gate 1A fails. Only item (1) is needed: it holds under N-q, and N-loc and N-free
are larger classes.
\end{enumerate}
The argument uses no Fisher-information or regularity assumption. It needs only that the
frozen class contains a free, label-independent, per-candidate distribution in the open simplex,
which is how specification \S5.3 describes it.
\end{proof}

\subsection{\texorpdfstring{What Lemma S1$'$ does \emph{not} claim}{What Lemma S1' does not claim}}
\begin{itemize}[leftmargin=*,itemsep=2pt]
\item It makes no claim under a \emph{structured} nuisance class. For example, if the circuit is
known and visibility is fixed at $1$, $\eps$ can be identified from absolute levels: the oracle
information is $0.24$--$0.62$ per shot at $\phi=0$. No structured deployable class was frozen,
and none is claimed blind.
\item Global non-identification is \emph{not} claimed for an arbitrary class that merely contains
a free $q_c$. That is where E006 narrows E001 \S2: under N-q (known directions), positivity can
bound the equivalence set, so only local non-identification is guaranteed.
\item It is not a statement about hardware, or about $\phi\neq0$.
\end{itemize}

\subsection{Numerical check (not a proof)}
Tool: \path{tools/audit_gate1a_efficient_information.py}; output: \path{results/errata/gate1a_efficient_information.json}.
At $\phi=0$, the largest efficient $\eps$ information over all frozen contexts and designs is
$6.8\times10^{-21}$ per shot (numerically zero). The oracle (known-response) information at the
same points is $0.24$--$0.62$ per shot.

%=====================================================================
\section{Lemma B (gap-shift bound)}

\subsection{Setting}
\begin{itemize}[leftmargin=*,itemsep=2pt]
\item \textbf{Circuits.} Candidate $j\in\{A,B\}$ is the frozen $p=1$ QAOA circuit on $|+\rangle^{\otimes n}$.
Its tagged angles $x_{j,1},\dots,x_{j,m_j}$ each enter through their own $R_Z(x_{j,i})$, inside a
compiled $\mathrm{CX}$--$R_Z$--$\mathrm{CX}$ block.
\item \textbf{Score.} The normalised cut score is $s(b)\in[0,1]$. The expected score and the expected gap are
\[
E_j(x_j)=\sum_b s(b)\,\big|\langle b|U_j(x_j)|{+}^{n}\rangle\big|^2\in[0,1],\qquad
g(x)=E_A(x_A)-E_B(x_B).
\]
\item \textbf{Error.} Target circuits carry no RC labels. Under the primary model, each tagged angle
$\theta_{j,i}>0$ becomes $\theta_{j,i}+\delta_{j,i}$ with
\[
\delta_{j,i}=\eps\,\theta_{j,i}+\phi\qquad\text{exactly (no linearisation).}
\]
\item \textbf{Notation.} $A:=\sum_{j,i}|\theta_{j,i}|$ and $M:=m_A+m_B$. The frozen threshold quantity is
$B(\eps,\phi):=|\eps|A+|\phi|M$.
\end{itemize}

\subsection{Statement}
\begin{lemmaB}
For all real $\eps,\phi$:
\[
\big|g(\theta+\delta)-g(\theta)\big| \;\le\; \tfrac12\,\norm{\delta}_1 \;\le\; \tfrac12\big(|\eps|A+|\phi|M\big) \;=\; \tfrac12 B(\eps,\phi).
\]
\end{lemmaB}

\subsection{Proof}
\begin{enumerate}[label=\textbf{Step \arabic*.},leftmargin=*,itemsep=4pt]
\item \textbf{Trigonometric degree in one angle.} Fix every argument except $x=x_{j,i}$. Then
$U_j=V_2\,R_Z(x)\,V_1$ with $R_Z(x)=\cos(x/2)\,I-i\sin(x/2)\,Z$. Each amplitude is
$\cos(x/2)\,\alpha_b-i\sin(x/2)\,\beta_b$ for some $\alpha_b,\beta_b$ independent of $x$. Its squared modulus is a linear
combination of $\cos^2(x/2)$, $\sin^2(x/2)$ and $\sin(x/2)\cos(x/2)$, so
\[
E_j = a+b\cos x+c\sin x = a + \rho\cos(x-x_0),\qquad \rho=\sqrt{b^2+c^2}.
\]
\item \textbf{Uniform derivative bound.} Since $0\le E_j\le1$ for \emph{all} $x$, the oscillation
$\max_x E_j-\min_x E_j = 2\rho \le 1$. Hence
\[
\Big|\tfrac{\partial E_j}{\partial x_{j,i}}\Big| \le \rho \le \tfrac12 \quad\text{at every point,}
\]
not only at $\theta$.
\item \textbf{Exact integral identity.} The two candidates depend on disjoint arguments, so every component of
$\nabla g(y)$ is $\pm$ one of the derivatives in Step 2, and $\norm{\nabla g(y)}_\infty\le\tfrac12$ for all $y$.
Along the segment $y(t)=\theta+t\delta$, $t\in[0,1]$, the fundamental theorem of calculus gives
\[
g(\theta+\delta)-g(\theta)=\int_0^1 \nabla g\big(y(t)\big)\cdot\delta\,dt,
\]
so, by H\"older's inequality,
\[
\big|g(\theta+\delta)-g(\theta)\big|\le \sup_y\norm{\nabla g(y)}_\infty\,\norm{\delta}_1 \le \tfrac12\norm{\delta}_1 .
\]
\item \textbf{Size of the perturbation.} $|\delta_{j,i}|\le|\eps|\,\theta_{j,i}+|\phi|$. Summing over all tagged rotations
gives $\norm{\delta}_1\le|\eps|A+|\phi|M$. \qed
\end{enumerate}

\subsection{Properties}
\begin{itemize}[leftmargin=*,itemsep=2pt]
\item \textbf{Norms:} $\ell_\infty$ on the gradient, $\ell_1$ on the perturbation (H\"older pairing).
\item \textbf{Constant:} $\tfrac12$ is \emph{uniform}. It is independent of the instance, the angles, $\eps$, $\phi$,
and the perturbation size, and depends only on the score range $[0,1]$.
\item \textbf{Remainder:} none. The bound is exact and \emph{global}; it is \emph{not} a first-order approximation.
\item \textbf{Consequence:} $B(\eps,\phi)$ is exactly twice a global bound at the \emph{true} parameters.
This corrects E003 \S6, which called it ``not a bound beyond first order''.
\end{itemize}

\subsection{Numerical check (not a proof)}
Tool: \path{tools/audit_gap_shift_bound.py}; output: \path{results/errata/gap_shift_bound_check.json}.
This is an exact evaluation of the frozen probability functions, with no sampling.

\medskip
\noindent\begin{tabular}{@{}lrr@{}}
\toprule
Population & Units with $(\eps,\phi)\neq0$ & $\max\,|\Delta g| \,/\, \big(B/2\big)$ \\
\midrule
Gate 2 & 6160 & 0.431 \\
Gate 3/4 & 4752 & 0.213 \\
\bottomrule
\end{tabular}

\medskip
The bound holds everywhere, with at least about $2.3\times$ slack in these populations.

\subsection{\texorpdfstring{What is \emph{not} proved about the frozen decision rule}{What is not proved about the frozen decision rule}}
\begin{enumerate}[leftmargin=*,itemsep=2pt]
\item $\hat B=B(\hat\eps,\hat\phi)$ uses \emph{estimates}. No coverage statement
$\Pr\big(B(\eps,\phi)\le\hat B\big)\ge1-\alpha$ has been proved for the frozen P-amp estimator.
\item The rule compares the \emph{sampled} initial gap $|\hat g_0|$ with $\hat B$. Lemma B concerns the
\emph{expected} gap. Nothing links that comparison to the probability of a decision flip, or to
correction benefit, under shot noise.
\item Lemma B says nothing about the residual error after recalibration with estimated parameters,
or about how tight the bound is.
\item It covers only the noiseless synthetic E2 model, with errors on tagged rotations only.
\end{enumerate}
\textbf{Correct description:} $\hat B$ is a \emph{heuristic plug-in alarm threshold built from a provable
population bound}, not a guarantee about the rule.

%=====================================================================
\newpage
\section{Review checklist}
\renewcommand{\arraystretch}{1.6}
\noindent\begin{tabularx}{\linewidth}{@{}c>{\raggedright\arraybackslash}Xc@{}}
\toprule
& Item to check & Initials / date \\
\midrule
$\square$ & S1$'$ setting: frame covariance gives label invariance at $\phi=0$. & \\
$\square$ & The three nuisance parameterizations faithfully cover the readings of specification \S5.3. & \\
$\square$ & (1) N-q construction: $\delta>0$; $q'_c$ in the open simplex; laws equal. & \\
$\square$ & (2) N-loc construction for $\eps'\neq0$ and for $\eps'=0$. & \\
$\square$ & (3) N-free construction uses only the single observed point per candidate. & \\
$\square$ & (4) The $\eps$ score is label independent and lies in the nuisance tangent span. & \\
$\square$ & Corollary: the nonunique profile MLE at the nine $\phi=0$ points implies FAIL under \S7 and \S9. & \\
$\square$ & ``Does not claim'' list accepted (structured classes; N-q only local). & \\
$\square$ & Lemma B Step 1: degree-1 trigonometric dependence on each tagged angle. & \\
$\square$ & Lemma B Step 2: oscillation $\le1$ implies $|\partial E|\le\tfrac12$ everywhere. & \\
$\square$ & Lemma B Step 3: disjoint arguments; the integral identity gives the global bound. & \\
$\square$ & Lemma B Step 4: $\delta_{j,i}=\eps\theta_{j,i}+\phi$ is exact in the target circuits (no labels). & \\
$\square$ & The unproved items about $\hat B$ are stated correctly and are not overclaimed elsewhere. & \\
\bottomrule
\end{tabularx}

\vspace{14pt}
\textbf{Decision:}\quad $\square$ accept as proved\quad $\square$ accept with corrections\quad $\square$ reject

\vspace{10pt}
\textbf{Notes:}

\end{document}
